\chapter{Metodologia}

\section{Classificação}

A pesquisa apresentada neste trabalho é classificada quanto à sua natureza, abordagem e fins:

\begin{itemize}
    \item \textbf{Quanto à natureza}, é uma pesquisa aplicada, pois busca desenvolver uma solução prática e direta para um problema real da Faculdade, visando otimizar o processo de gerenciamento de estágios e suas diversas etapas.
    
    \item \textbf{Quanto à abordagem}, é uma pesquisa qualitativa, com o objetivo de compreender de forma detalhada as necessidades dos stakeholders (alunos, professores e coordenadores) no gerenciamento de estágios, coletando dados por meio de entrevistas, questionários e análise de processos administrativos.
    
    \item \textbf{Quanto aos fins}, é exploratória, pois busca investigar a viabilidade e os desafios de implementação de um sistema de gerenciamento de estágios, analisando as melhores práticas, tecnologias e processos administrativos, com foco na melhoria da gestão e acompanhamento das atividades de estágio.
\end{itemize}

O desenvolvimento dessa ferramenta visa facilitar o gerenciamento das informações relacionadas aos estágios, proporcionando uma solução centralizada e eficiente para todos os envolvidos no processo, desde a captação das vagas até o acompanhamento das atividades dos estagiários.

\section{Metodologia para Pesquisa do Referencial Teórico}


Esta fase construtiva envolve a elaboração de um plano de pesquisa e a realização efetiva da pesquisa, sendo desenvolvida em duas etapas distintas. A exploração dos conteúdos nessas etapas foi conduzida por meio de pesquisa em livros, artigos acadêmicos, sites governamentais e outros recursos relevantes que abordam temas relacionados ao gerenciamento de estágios, assinatura digital e Login Único.

A primeira etapa teve como objetivo compreender a implementação do \textit{Login Único} (SSO) no contexto da autenticação de usuários no sistema de gerenciamento de estágios. A pesquisa envolveu a análise de como o Login Único pode ser utilizado para simplificar o processo de login, e possibilitando a utilizaçao de outra APIs do GOV como a assinatura digital integrando o sistema com outra plataformas, oferecendo um acesso mais rápido e seguro para os usuários.

Na segunda etapa, foi realizado um estudo sobre o funcionamento das assinaturas digitais e o Login Único, com foco no uso da \textit{Assinatura Digital do GOV}. O objetivo era entender como a tecnologia de assinatura digital pode ser aplicada na validação de documentos relacionados aos estágios (como contratos, relatórios e formulários), garantindo a autenticidade e a integridade das informações, bem como a conformidade com a legislação brasileira. Este estudo envolveu a análise de como a assinatura digital pode ser integrada ao sistema de gerenciamento de estágios e como ela facilita a validação de documentos no ambiente acadêmico e corporativo.

Essa abordagem metodológica permitiu não apenas a compreensão dos aspectos técnicos envolvidos na implementação da assinatura digital e Login Único, mas também a consideração dos aspectos legais e regulatórios relacionados ao uso dessas tecnologias no contexto acadêmico, assegurando a conformidade do sistema de gerenciamento de estágios com as normas e leis vigentes no Brasil.


\section{Metodologia de Desenvolvimento}

A metodologia de desenvolvimento adotada para este projeto é o \textit{Scrum}, uma abordagem ágil que visa proporcionar entregas rápidas e eficientes, focando em ciclos de desenvolvimento curtos e na colaboração constante entre a equipe e os stakeholders. O Scrum é uma metodologia que se baseia em princípios como transparência, inspeção e adaptação, com o objetivo de entregar software de alta qualidade de forma incremental, adaptando-se às mudanças nas necessidades ao longo do processo.

\subsection{Scrum: Visão Geral}

O \textit{Scrum} é um framework ágil que organiza o trabalho em ciclos curtos chamados de \textit{sprints}, com duração geralmente de 2 a 4 semanas. Cada sprint resulta em uma entrega funcional do sistema, permitindo que os stakeholders possam avaliar o progresso e fornecer feedback. O processo Scrum é estruturado em torno de papéis, eventos e artefatos, que ajudam a organizar e monitorar o desenvolvimento.

Os principais papéis no Scrum incluem:

\begin{itemize}
    \item \textbf{Product Owner (PO)}: Responsável por gerenciar o backlog do produto, priorizando as funcionalidades e garantindo que a equipe de desenvolvimento esteja focada nas tarefas mais importantes para o negócio.
    \item \textbf{Scrum Master}: Facilita o processo Scrum, removendo obstáculos para a equipe e garantindo que as práticas ágeis sejam seguidas.
    \item \textbf{Time de Desenvolvimento}: Equipe multifuncional que trabalha para entregar as funcionalidades solicitadas, composta por desenvolvedores, testadores, e outros profissionais necessários.
\end{itemize}

Os principais eventos no Scrum incluem:

\begin{itemize}
    \item \textbf{Sprint Planning}: Reunião no início de cada sprint, onde as tarefas são definidas e o trabalho a ser realizado é planejado.
    \item \textbf{Daily Scrum}: Reunião diária de 15 minutos, onde a equipe compartilha o progresso, os obstáculos encontrados e o que será feito no dia.
    \item \textbf{Sprint Review}: Reunião ao final de cada sprint para demonstrar as funcionalidades entregues e coletar feedback dos stakeholders.
    \item \textbf{Sprint Retrospective}: Reunião de reflexão após cada sprint, onde a equipe discute o que funcionou bem, o que pode ser melhorado e ações para o próximo sprint.
\end{itemize}


\subsection{Aplicação do Scrum}

No contexto do desenvolvimento do software de gerenciamento de estágios, o Scrum será utilizado  de maneira adaptada e reduzida para garantir que o processo de desenvolvimento seja eficiente, flexível e alinhado às necessidades da FCTE e dos alunos. A equipe de desenvolvimento trabalhará em sprints, realizando entregas incrementais de funcionalidades, com o objetivo de entregar uma versão funcional do sistema ao final de cada ciclo.

O \textit{Product Owner} será responsável por gerenciar o backlog do produto, priorizando funcionalidades essenciais para o gerenciamento eficiente de estágios, como o cadastro de alunos, empresas, e a geração de documentos relacionados aos estágios.

A cada sprint, a equipe de desenvolvimento trabalhará na implementação de funcionalidades específicas do sistema, como integração com a assinatura digital do GOV, integração com o Login Único, e gerenciamento das vagas de estágio. Durante as \textit{Daily Scrums}, a equipe discutirá o progresso, identificará desafios e ajustará o plano conforme necessário.
